% GENERATED FILE. Edit canonical lessons, then run npm run generate:books.
% canonical-source-hash: fnv1a64:18ac19619363f91c
% canonical-lessons: TE-C543-kaceri, TE-C543-cappatlu, TE-C543-abhinandincu, TE-C543-asvadincu, TE-C543-kolatam

\chapter{At a Performance}
\label{ch:te-at-a-performance}
\hlchaptermodality{543}

\begin{chapteropening}
\textbf{By the end of this chapter:} \emph{I can talk about a concert, the applause and the folk dances.}

Complete the last lesson of chapter 543: I can talk about a concert, the applause and the folk dances.
\end{chapteropening}

\section[\texorpdfstring{\te{కచేరీ} (kacērī)}{కచేరీ (kacērī)}]{\te{కచేరీ} (kacērī) — a concert; a recital}

\label{lesson:TE-C543-kaceri}



\begin{quote}
Before the new one: say the Telugu for a sculptor, then the Telugu for a masterpiece.
\end{quote}

\subsection*{You'll want to know: \te{కచేరీ}}
\textbf{\te{కచేరీ}} — \emph{kacērī} — \textquotedblleft{}a concert; a recital\textquotedblright{}.

\begin{quote}
\te{శనివారం} \te{సాయంత్రం} \te{సంగీత} \te{కచేరీ} \te{జరిగింది}.
\end{quote}

\begin{quote}
\emph{śanivāraṁ sāyantraṁ saṅgīta kacērī jarigindi.}
\end{quote}

\begin{quote}
There was a music concert on Saturday evening.
\end{quote}

Also an old word for a government office.

\begin{cousinweb}
\textbf{\te{కచేరీ}} (\emph{kacērī}) came into Telugu through Hindustani \emph{kacahrī}, \textquotedblleft{}a court, an office\textquotedblright{}; its earlier origin is not settled.
\end{cousinweb}

\begin{grammarlens}[title={What you've built}]
One more word for describing what you lived and felt.
\end{grammarlens}

\subsection*{Your turn}
\begin{itemize}
\raggedright
  \item \emph{Say it:} \emph{kacērī}
  \item \emph{Say it:} \emph{kacērī}, once more
  \item \emph{Say it:} \emph{śanivāraṁ sāyantraṁ saṅgīta kacērī jarigindi.} — There was a music concert on Saturday evening.
  \item \emph{Recall:} say the Telugu for a sculptor, then the Telugu for a masterpiece, then say \emph{kacērī} again
\end{itemize}

\begin{tcolorbox}[breakable,colback=teal!4,colframe=teal!35!black,title={Before you move on}]
What does \te{కచేరీ} mean? (\textquotedblleft{}A concert; a recital\textquotedblright{}.) Say it once more.
\end{tcolorbox}
\section[\texorpdfstring{\te{చప్పట్లు} (cappaṭlu)}{చప్పట్లు (cappaṭlu)}]{\te{చప్పట్లు} (cappaṭlu) — applause; clapping}

\label{lesson:TE-C543-cappatlu}



\begin{quote}
Before the new one: say the Telugu for a masterpiece, then the Telugu for a concert.
\end{quote}

\subsection*{You'll want to know: \te{చప్పట్లు}}
\textbf{\te{చప్పట్లు}} — \emph{cappaṭlu} — \textquotedblleft{}applause; clapping\textquotedblright{}.

\begin{quote}
\te{పాట} \te{అయ్యాక} \te{అందరూ} \te{చప్పట్లు} \te{కొట్టారు}.
\end{quote}

\begin{quote}
\emph{pāṭa ayyāka andarū cappaṭlu koṭṭāru.}
\end{quote}

\begin{quote}
When the song ended, everyone clapped.
\end{quote}

\te{చప్పట్లు} \te{కొట్టు}, to clap.

\begin{cousinweb}
\textbf{\te{చప్పట్లు}} (\emph{cappaṭlu}) is native Dravidian: DEDR 2335 lists Telugu \emph{cappaṭlu}, \textquotedblleft{}clapping of the hands\textquotedblright{}, with Tamil \emph{cappāṇi}.
\end{cousinweb}

\begin{grammarlens}[title={What you've built}]
One more word for describing what you lived and felt.
\end{grammarlens}

\subsection*{Your turn}
\begin{itemize}
\raggedright
  \item \emph{Say it:} \emph{cappaṭlu}
  \item \emph{Say it:} \emph{cappaṭlu}, once more
  \item \emph{Say it:} \emph{pāṭa ayyāka andarū cappaṭlu koṭṭāru.} — When the song ended, everyone clapped.
  \item \emph{Recall:} say the Telugu for a masterpiece, then the Telugu for a concert, then say \emph{cappaṭlu} again
\end{itemize}

\begin{tcolorbox}[breakable,colback=teal!4,colframe=teal!35!black,title={Before you move on}]
What does \te{చప్పట్లు} mean? (\textquotedblleft{}Applause; clapping\textquotedblright{}.) Say it once more.
\end{tcolorbox}
\section[\texorpdfstring{\te{అభినందించు} (abhinandiñcu)}{అభినందించు (abhinandiñcu)}]{\te{అభినందించు} (abhinandiñcu) — to congratulate}

\label{lesson:TE-C543-abhinandincu}



\begin{quote}
Before the new one: say the Telugu for a concert, then the Telugu for applause.
\end{quote}

\subsection*{You'll want to know: \te{అభినందించు}}
\textbf{\te{అభినందించు}} — \emph{abhinandiñcu} — \textquotedblleft{}to congratulate\textquotedblright{}.

\begin{quote}
\te{గెలిచిన} \te{జట్టును} \te{అందరూ} \te{అభినందించారు}.
\end{quote}

\begin{quote}
\emph{gelicina jaṭṭunu andarū abhinandiñcāru.}
\end{quote}

\begin{quote}
Everyone congratulated the winning team.
\end{quote}

Past (I …): \te{అభినందించాను}, \emph{abhinandiñcānu}.

\te{అభినందనలు}, congratulations.

\begin{cousinweb}
\textbf{\te{అభినందించు}} (\emph{abhinandiñcu}) is built on Sanskrit \emph{abhinandana}, \textquotedblleft{}delighting; praising, applauding\textquotedblright{} (Monier-Williams), with the Telugu verb ending -\te{ఇంచు}.
\end{cousinweb}

\begin{grammarlens}[title={What you've built}]
One more word for describing what you lived and felt.
\end{grammarlens}

\subsection*{Your turn}
\begin{itemize}
\raggedright
  \item \emph{Say it:} \emph{abhinandiñcu}
  \item \emph{Say it:} \emph{abhinandiñcu}, once more
  \item \emph{Say it:} \emph{gelicina jaṭṭunu andarū abhinandiñcāru.} — Everyone congratulated the winning team.
  \item \emph{Recall:} say the Telugu for a concert, then the Telugu for applause, then say \emph{abhinandiñcu} again
\end{itemize}

\begin{tcolorbox}[breakable,colback=teal!4,colframe=teal!35!black,title={Before you move on}]
What does \te{అభినందించు} mean? (\textquotedblleft{}To congratulate\textquotedblright{}.) Say it once more.
\end{tcolorbox}
\section[\texorpdfstring{\te{ఆస్వాదించు} (āsvādiñcu)}{ఆస్వాదించు (āsvādiñcu)}]{\te{ఆస్వాదించు} (āsvādiñcu) — to savour; to relish}

\label{lesson:TE-C543-asvadincu}



\begin{quote}
Before the new one: say the Telugu for applause, then the Telugu for to congratulate.
\end{quote}

\subsection*{You'll want to know: \te{ఆస్వాదించు}}
\textbf{\te{ఆస్వాదించు}} — \emph{āsvādiñcu} — \textquotedblleft{}to savour; to relish\textquotedblright{}.

\begin{quote}
\te{మేము} \te{ఆ} \te{సంగీతాన్ని} \te{హాయిగా} \te{ఆస్వాదించాము}.
\end{quote}

\begin{quote}
\emph{mēmu ā saṅgītānni hāyigā āsvādiñcāmu.}
\end{quote}

\begin{quote}
We relished that music at our ease.
\end{quote}

Past (I …): \te{ఆస్వాదించాను}, \emph{āsvādiñcānu}.

Enjoying food, music or a view slowly.

\begin{cousinweb}
\textbf{\te{ఆస్వాదించు}} (\emph{āsvādiñcu}) is built on Sanskrit \emph{āsvādana}, \textquotedblleft{}eating, tasting, enjoying\textquotedblright{} (Monier-Williams), with the Telugu verb ending -\te{ఇంచు}.
\end{cousinweb}

\begin{grammarlens}[title={What you've built}]
One more word for describing what you lived and felt.
\end{grammarlens}

\subsection*{Your turn}
\begin{itemize}
\raggedright
  \item \emph{Say it:} \emph{āsvādiñcu}
  \item \emph{Say it:} \emph{āsvādiñcu}, once more
  \item \emph{Say it:} \emph{mēmu ā saṅgītānni hāyigā āsvādiñcāmu.} — We relished that music at our ease.
  \item \emph{Recall:} say the Telugu for applause, then the Telugu for to congratulate, then say \emph{āsvādiñcu} again
\end{itemize}

\begin{tcolorbox}[breakable,colback=teal!4,colframe=teal!35!black,title={Before you move on}]
What does \te{ఆస్వాదించు} mean? (\textquotedblleft{}To savour; to relish\textquotedblright{}.) Say it once more.
\end{tcolorbox}
\section[\texorpdfstring{\te{కోలాటం} (kōlāṭaṁ)}{కోలాటం (kōlāṭaṁ)}]{\te{కోలాటం} (kōlāṭaṁ) — a stick dance}

\label{lesson:TE-C543-kolatam}



\begin{quote}
Before the new one: say the Telugu for to congratulate, then the Telugu for to savour.
\end{quote}

\subsection*{You'll want to know: \te{కోలాటం}}
\textbf{\te{కోలాటం}} — \emph{kōlāṭaṁ} — \textquotedblleft{}a stick dance\textquotedblright{}.

\begin{quote}
\te{పండుగలో} \te{ఆడపిల్లలు} \te{కోలాటం} \te{ఆడారు}.
\end{quote}

\begin{quote}
\emph{paṇḍugalō āḍapillalu kōlāṭaṁ āḍāru.}
\end{quote}

\begin{quote}
At the festival the girls did the stick dance.
\end{quote}

Dancers strike short sticks in rhythm: \te{కోలు}, stick, + \te{ఆటం}, play.

\begin{cousinweb}
\textbf{\te{కోలాటం}} (\emph{kōlāṭaṁ}) is built from \te{కోలు}, \textquotedblleft{}a stick\textquotedblright{}, and native \emph{āṭaṁ}, \textquotedblleft{}play, dance\textquotedblright{} (DEDR 347).
\end{cousinweb}

\begin{grammarlens}[title={What you've built}]
One more word for describing what you lived and felt.
\end{grammarlens}

\subsection*{Your turn}
\begin{itemize}
\raggedright
  \item \emph{Say it:} \emph{kōlāṭaṁ}
  \item \emph{Say it:} \emph{kōlāṭaṁ}, once more
  \item \emph{Say it:} \emph{paṇḍugalō āḍapillalu kōlāṭaṁ āḍāru.} — At the festival the girls did the stick dance.
  \item \emph{Recall:} say the Telugu for to congratulate, then the Telugu for to savour, then say \emph{kōlāṭaṁ} again
\end{itemize}

\begin{tcolorbox}[breakable,colback=teal!4,colframe=teal!35!black,title={Before you move on}]
What does \te{కోలాటం} mean? (\textquotedblleft{}A stick dance\textquotedblright{}.) Say it once more.
\end{tcolorbox}
